\documentclass[10pt]{beamer}

\usepackage[latin1]{inputenc}
\usepackage[spanish]{babel}

\mode<presentation>
\usetheme{Madrid}
%\usecolortheme[RGB={111,73,135}]{structure}
\usecolortheme[RGB={128,0,0}]{structure}
%\usecolortheme[RGB={0,96,0}]{structure}
%\usecolortheme[RGB={200,0,200}]{structure}
%\usecolortheme[RGB={0,128,0}]{structure}
%\usecolortheme[RGB={0,0,128}]{structure}
\usefonttheme{serif}
\useinnertheme{rectangles}
\useoutertheme{split}
\setbeamercovered{transparent}


% Definiciones para usar luego
\newtheorem{ejemplo}{Ejemplo}
\newtheorem{definicion}{Definici�n}

\title{An�lisis de redes con WireShark}
\author{Jes�s Espino Garc�a}
\date{10 de Marzo de 2009}
\subject{An�lisis de redes con WireShark}

\institute[GUL UC3M]{
  Grupo de usuarios de Linux\\
  Universidad Carlos III de Madrid.\\
  \ \\
  \includegraphics[height=1.5cm]{gul}}

\begin{document}

  \frame{\maketitle}
  
  \section{Introducci�n}

  \begin{frame}{Conceptos b�sicos de red}
    \begin{itemize}
      \item Pila TCP-UDP/IP (F�sica,Enlace,IP,TCP,Aplicaci�n)
      \item Hardware (Hub/Switch)
      \item Puerto.
      \item Protocolo.
    \end{itemize}
  \end{frame}

  \begin{frame}{�Qu� es?}
    \begin{itemize}
      \item Sniffer de red.
      \item Es GPL.
      \item Multiplataforma.
      \item Reconoce mas de 480 protocolos.
      \item Capaz de seguir sesiones TCP.
      \item Soporta archivos de captura de otros sniffers.
      \item Compatible con tcpdump.
    \end{itemize}
  \end{frame}

  \begin{frame}{�Por qu�?}
    \begin{itemize}
      \item Detecci�n de errores.
      \item Detecci�n de intrusiones.
      \item Detecci�n de congestiones.
      \item Obtenci�n de informaci�n privada.
    \end{itemize}
  \end{frame}

  \begin{frame}{�C�mo?}
    \begin{itemize}
      \item Usando wireshark.
      \item Usando tcpdump.
    \end{itemize}
  \end{frame}

  \begin{frame}{�D�nde?}
    \begin{itemize}
      \item Misma maquina.
      \item Misma red f�sica (hubs).
      \item Mismo switch con puerto espejo.
      \item Mismo switch con cache poisoning.
      \item En ambas maquinas (origen y destino).
    \end{itemize}
  \end{frame}

  \section{Obteniendo el trafico de red}

  \begin{frame}[containsverbatim]{Obteniendo el trafico de red}
    \begin{block}{Con wireshark}
      \begin{itemize}
        \item Abrir WireShark como root.
        \item Ir a Capture > Interface\dots
        \item Hacer click en start en la interfaz que nos interesa.
        \item Al terminar ir a Capture > Stop.
      \end{itemize}
    \end{block}

    \begin{block}{Con tcpdump}
      \begin{itemize}
        \item En un terminal como root ejecutar \verb+tcpdump -w <fichero.pcap>+
        \item Abrir el fichero creado con WireShark.
      \end{itemize}
    \end{block}
  \end{frame}
  
  \section{Usando WireShark}

  \subsection{Filtrando paquetes}

  \begin{frame}{Filtrando paquetes}
    \begin{block}{En la captura}
      \begin{itemize}
        \item Filtros que permiten capturar solo el trafico que nos interesa.
        \item Generan ficheros mas peque�os y manejables.
        \item Pueden perderse tramas interesantes para el diagnostico.
      \end{itemize}
    \end{block}

    \begin{block}{En el mostrado}
      \begin{itemize}
        \item Ocultan los paquetes que no coinciden con el filtro.
        \item Maneja todos los datos, podemos diagnosticar con mas informaci�n.
        \item Los ficheros normalmente son mas grandes y mas pesados de manejar.
      \end{itemize}
    \end{block}
  \end{frame}

  \begin{frame}{Filtrando paquetes}
    \begin{block}{Expresiones}
      \begin{itemize}
        \item Para poder filtrar mas f�cilmente podemos usar el bot�n "Expresion\dots".
        \item Nos permite construir filtros de manera sencilla sin tener que acordarnos de todo.
      \end{itemize}
    \end{block}

    \begin{block}{Seguir tramas}
      \begin{itemize}
        \item Hacemos click derecho sobre un paquete y seleccionamos "Follow TCP Stream".
        \item Nos crea autom�ticamente un filtro para esta trama.
        \item Ademas nos abre una ventana con la informaci�n intercambiada.
      \end{itemize}
    \end{block}
  \end{frame}
  
  \subsection{Modificando los tiempos}

  \begin{frame}{Modificando los tiempos}
    \begin{block}{El formato}
      \begin{itemize}
        \item Se pueden seleccionar diferentes formatos y precisiones para los tiempos.
        \item Formatos como, la fecha, la hora, el numero de segundos desde el principio de la captura\dots
        \item Precisiones desde segundos hasta nanosegundos.
      \end{itemize}
    \end{block}

    \begin{block}{La referencia}
      \begin{itemize}
        \item Tambi�n es interesante contar el tiempo a partir de un determinado paquete.
        \item WireShark permite seleccionar un paquete para que sea el segundo 0.
      \end{itemize}
    \end{block}
  \end{frame}

  \section{Ejemplos}

  \begin{frame}{Ejemplos}
    \begin{itemize}
      \item Telnet.
      \item IMAP.
      \item DNS.
      \item DHCP.
      \item Virus Slammer.
      \item WakeOnLan.
    \end{itemize}
  \end{frame}

  \section{Para terminar.}

  \begin{frame}[containsverbatim]{Referencias}
    \begin{itemize}
      \item \verb+man wireshark+: Pagina del manual de WireShark.
      \item \verb+man tcpdump+: Pagina del manual de tcpdump.
      \item http://www.wireshark.org: Pagina oficial de WireShark.
      \item http://wiki.wireshark.org/SampleCaptures: Ejemplos de capturas.
    \end{itemize}
  \end{frame}

  \begin{frame}{Dudas}
    \dots
  \end{frame}
  
  \begin{frame}
    \frametitle{\begin{center}Fin\end{center}}
    \begin{center}
      \includegraphics[height=3cm]{gul}
    \end{center}
  \end{frame}

\end{document}
